% =====================================================================
% resumen.tex
% Responsabilidad única: resumen ejecutivo (150 a 250 palabras, como pide
% el enunciado): problema, datos, modelos, resultado, recomendación y
% limitación principal. Todas las cifras son macros.
% =====================================================================
\begin{abstract}
\noindent
\textbf{Resumen ejecutivo.} Un laboratorio de microbiología necesita saber
qué muestras van a resultar no conformes para revisarlas primero y cuántas
muestras van a llegar para programar su trabajo. Se trabajó con la
exportación real anonimizada del laboratorio y, como complemento para aplicar
más técnicas del curso, con una base simulada de la misma estructura que
agrega los controles de calidad y los datos de recepción que la exportación
no trae (cribado de \CriCandidatos{} análisis con una regla fijada de
antemano, arreglos L16 de Taguchi, Bayes, validación por bloques temporales y
una sola evaluación en prueba). En la base real (\RealMuestras{} muestras de
\RealClientes{} clientes en el periodo de desarrollo) las muestras no
llegan como una Poisson pero los lotes sí así que la llegada es una Poisson
compuesta. Para clasificar la no conformidad al ingreso se compararon una
referencia, la regla de la tasa histórica del cliente, regresiones logísticas
y árboles. La logística congelada ordenó el riesgo de 2026 con un ROC-AUC de
\RealPruebaAUC{} y una PR-AUC de \RealPruebaAP{}, frente a una prevalencia de
\RealPruebaPrevalencia{}, pero no se distingue de la regla simple
(\RealPruebaAPRegla{}) y su umbral dejó de pagar en 2026. Se recomienda usar
la puntuación para ordenar la revisión y no como alarma automática. La
limitación principal es la deriva de 2026, con menos no conformidad y más
clientes nuevos, que obliga a recalibrar antes de cualquier uso.
\end{abstract}
